\documentclass[10pt]{article}
\usepackage[russian]{babel}
\usepackage[utf8]{inputenc}
\usepackage[T2A]{fontenc}
\usepackage{amsmath}
\usepackage{amsfonts}
\usepackage{amssymb}
\usepackage[version=4]{mhchem}
\usepackage{stmaryrd}
\usepackage{bbold}
\usepackage{enumerate}

\begin{document}

\section{Задача 1 (3 балла)}

\subsection{Пункт 1.1 (1 балл)}
\subsubsection{Условие}

\[
S=\left\{x \in \mathbb{R}^{d} \mid x_{1} \leq 2 x_{2} \leq \ldots \leq d x_{d}\right\}
\]

\vspace{5mm}
\subsubsection{Решение}

пусть $x, y \in S$, тогда докажем что $\alpha x + (1-\alpha)y \in S$

по условию $kx_k < (k+1)x_{k+1}, ky_k < (k+1)y_{k+1}$

очевидно что тогда $\alpha kx_k + (1-\alpha)ky_k < \alpha (k+1)x_{k+1} + (1-\alpha) (k+1)y_{k+1}$

откуда и следует что $\alpha x + (1-\alpha)y \in S$

\subsection{Пункт 1.2 (1 балл)}
\subsubsection{Условие}

\[
S=\left\{x \in \mathbb{R}^{2} \mid\left(x_{1}-x_{2}\right)^{2}+\left(x_{2}-1\right)^{2} \leq 1\right\}
\]

\subsubsection{Решение}

применим преобразование к S
\[
\begin{pmatrix} y_1 \\ y_2 \end{pmatrix} = 
\begin{pmatrix} 1 & -1 \\ 0 & 1 \end{pmatrix}
\begin{pmatrix} x_1 \\ x_2 \end{pmatrix} + 
\begin{pmatrix} 0 \\ -1 \end{pmatrix}.
\]

тогда такое преобразование переводит наше множество в единичный круг (что и является выпуклым множеством). Тогда S - является выпууклым множеством как прообраз аффинного преобразования от выпуклого множества

\subsection{Пункт 1.3 (1 балл)}
\subsubsection{Условие}

\[
S=\left\{x \in \mathbb{R}^{2} \mid x_{1}^{2} \leq x_{2} \leq x_{1}+4\right\}
\]

\subsubsection{Решение}

рассмотрим два множества

\[
S_1=\left\{x \in \mathbb{R}^{2} \mid x_{1}^{2} \leq x_{2} \right\}
\]

\[
S_2=\left\{x \in \mathbb{R}^{2} \mid x_{2} \leq x_{1}+4\right\}
\]

$S_1$ - Эпиграф $x_2 = x_1^2$ (выпуклой функции). Значит $S_1$ - выпукло

$S_2$ - полуплоскость. $S_2$ тоже выпукло

так как $S$ = $S_1 \cap S_2$ то $S$ тоже выпукло




\section{Задача 2 (2 балла)}
Пусть $S \subseteq \mathbb{R}^{d}$ и $\|\cdot\|$ — произвольная норма на $\mathbb{R}^{d}$.

\subsection{Пункт 2.1 (1 балл)}
\subsubsection{Условие}

Для $a \geq 0$ определим множество $S_{a}$ как:
\[
S_{a}=\{x \mid \operatorname{dist}(x, S) \leq a\}, \text{ где } \operatorname{dist}(x, S)=\inf_{y \in S}\|x-y\|
\]
Множество $S_{a}$ называется расширенным на $a$ относительно $S$. Докажите, что если $S$ выпукло, то $S_{a}$ также выпукло.

\subsubsection{Решение}

Тогда по определению инфимума

$\forall \varepsilon $ существует $ x_{\varepsilon}$ : $| x - x_{\varepsilon}| < a + \varepsilon$

$\forall \varepsilon $ существует $ y_{\varepsilon}$ : $| y - y_{\varepsilon}| < a + \varepsilon $

пусть $x, y \in S_a$. Докажем что $\alpha x + (1 - \alpha)y \in S_a $

зафиксируем $\varepsilon$

для $x, y$ существуют соответсвующие $x_{\varepsilon}, y_{\varepsilon}$

$$

| \alpha x + (1 - \alpha)y - \alpha x_{\varepsilon} - (1 - \alpha)y_{\varepsilon} | <

| \alpha (x - x_{\varepsilon}) + (1 - \alpha)(y - y_{\varepsilon})| <

\alpha| x - x_{\varepsilon} | + (1 - \alpha)|y - y_{\varepsilon}| < 

\alpha (a + \varepsilon) + (1  -\alpha)(a + \varepsilon) = 

(a + \varepsilon)
$$

по итогу мы доказали что 

$\forall \varepsilon $ существует $t_{\varepsilon}$ : $| \alpha x + (1 - \alpha)y - t_{\varepsilon}| < a + \varepsilon $

где $t_{\varepsilon} =  \alpha x_{\varepsilon} + (1 - \alpha)y_{\varepsilon}$

получается $ \alpha x + (1 - \alpha)y \in S_a$


\subsection{Пункт 2.2 (1 балл)}
\subsubsection{Условие}

Для $a \geq 0$ определим множество $S_{-a}$ как:
\[
S_{-a}=\{x \mid B(x, a) \subset S\}
\]
где $B(x, a)$ — открытый шар. Множество $S_{-a}$ называется суженным на $a$ относительно $S$. Докажите, что если $S$ выпукло, то $S_{-a}$ также выпукло.

\subsubsection{Решение}

Возьмем произвольные точки $x_1, x_2 \in S_{-a}$. По определению:
$$
B(x_1, a) \subset S, B(x_2, a) \subset S.
$$

$x_t = t x_1 + (1 - t) x_2$. Покажем, что $x_t \in S_{-a}$, т.е. $B(x_t, a) \subset S$.

Возьмем произвольную точку $z \in B(x_t, a)$. Тогда: $|z - x_t| < a.$

Заметим, что $z$ можно представить в виде:

$z = x_t + v, \quad \text{где } \|v\| < a.$

Тогда:

$z = t x_1 + (1 - t) x_2 + v$

$z = t(x_1 + v) + (1 - t)(x_2 + v)$


Так как $x_1, x_2 \in S_{-a}$, то для любого $v$ с $|v| < a$ выполняется:

$x_1 + v \in B(x_1, a) \subset S$

$x_2 + v \in B(x_2, a) \subset S.$


Поскольку $S$ выпукло, $t(x_1 + v) + (1 - t)(x_2 + v)$ также принадлежит $S$, 
т.е. $z \in S$.

Таким образом, для любого $z \in B(x_t, a)$ выполнено $z \in S$, значит:

$B(x_t, a) \subset S.$

Следовательно, $x_t \in S_{-a}$, что доказывает выпуклость $S_{-a}$.

\section{Задача 3 (1 балл)}

\subsubsection{Условие}

Пусть $X \subseteq \mathbb{R}^{d}$ и $x^{0} \in X$. Докажите, что множество
\[
K\left(X, x^{0}\right)=\left\{y \in \mathbb{R}^{d} \mid y^{\top} x^{0} \geq y^{\top} x \ \forall x \in X\right\}
\]
является выпуклым конусом.

\subsubsection{Решение}

Докажем, что:

Для любого $y \in K(X, x^0)$ и $\lambda \geq 0$ выполняется $\lambda y \in K(X, x^0)$ (свойство конуса).

Для любых $y_1, y_2 \in K(X, x^0)$ и $\lambda_1, \lambda_2 \geq 0$ выполняется $\lambda_1 y_1 + \lambda_2 y_2 \in K(X, x^0)$ (свойство выпуклого конуса).


1(свойство конуса)

Пусть $y \in K(X, x^0)$ и $\lambda \geq 0$. По определению $K(X, x^0)$ для любого $x \in X$:

$y^T x^0 \geq y^T x.$

Умножая обе части на $\lambda \geq 0$, получаем:

$\lambda y^T x^0 \geq \lambda y^T x,$

что эквивалентно:

$(\lambda y)^\top x^0 \geq (\lambda y)^\top x \quad \forall x \in X.$

Следовательно, $\lambda y \in K(X, x^0)$.

2(свойство выпуклого конуса)

Пусть $y_1, y_2 \in K(X, x^0)$ и $\lambda_1, \lambda_2 \geq 0$. Рассмотрим произвольный $x \in X$. По определению $K(X, x^0)$:

$y_1^T x^0 \geq y_1^T x \quad, \quad y_2^T x^0 \geq y_2^T x.$

Умножая первое неравенство на $\lambda_1 \geq 0$, а второе на $\lambda_2 \geq 0$, получаем:

$\lambda_1 y_1^T x^0 \geq \lambda_1 y_1^T x \quad \text{и} \quad \lambda_2 y_2^T x^0 \geq \lambda_2 y_2^T x.$

Складывая эти неравенства:

$\lambda_1 y_1^T x^0 + \lambda_2 y_2^T x^0 \geq \lambda_1 y_1^T x + \lambda_2 y_2^T x.$

Это эквивалентно:
$$
(\lambda_1 y_1 + \lambda_2 y_2)^T x^0 \geq (\lambda_1 y_1 + \lambda_2 y_2)^T x.
$$
Так как $x \in X$ произвольный, то $\lambda_1 y_1 + \lambda_2 y_2 \in K(X, x^0)$.

Таким образом, $K(X, x^0)$ является выпуклым конусом.

\section{Задача 4 (2 балла)}

\subsubsection{Условие}

Назовём множество $X \subseteq \mathbb{R}^{d}$ «средневыпуклым», если для любых его элементов $x$ и $y$ их середина также принадлежит $X$, то есть $\frac{x+y}{2} \in X$. Докажите, что для замкнутых множеств «средневыпуклость» равносильна выпуклости.

\subsubsection{Решение}

пусть $x, y \in X$. Покажем что $t = \alpha x + (1-\alpha) \in X$

рассмотрим такую последовательность

$x_0 = \frac{x+y}{2}$
left = x
right = y
если$ t \in [left, x_0]$
тогда $right = x_0$
иначе $left = x_0$

далее \forall n натурального
x_n = \frac{left + right}{2}
если$ t \in [left, x_{n}]$
тогда $right = x_n$
иначе $left = x_n$

посути мы генерируем последовательность бинарным поиском котора ясходится к $t$

по своиству множества все лементы принадлежат $X$. Тогда мы нашли $\{ x_n \}_{n=0} ё\subset X$ которая сходится к  $\alpha x + (1-\alpha)$

в силу замкнутости $X$ $\alpha x + (1-\alpha) \in X$
тем самым мы доказали выпуклость множества

\section{Задача 5 (2 балла)}

\subsubsection{Условие}

Пусть $X \in \mathbb{R}^{n \times d}, y \in \mathbb{R}^{n}$. Докажите, что будет выполнено ровно одно из двух условий:
\begin{enumerate}
  \item Найдётся $v \in \mathbb{R}_{+}^{d}$, что $X v=y$,
  \item Найдётся $u \in \mathbb{R}^{n}$, что $u^{\top} X \in \mathbb{R}_{+}^{d}$ и $\langle u, y\rangle<0$.
\end{enumerate}
Указание: используйте теорему о строгой отделимости из пособия.

\subsubsection{Решение}

Рассмотрим  $K = \{Xv \mid v \in \mathbb{R}_+^d\}.$

по условию $y \in K$.

Если условие (1) не выполнено, то $y \notin K$. По теореме о строгой отделимости существует $u \in R^n$ такой, что $u^T z \geq 0 \forall z \in K, u^T y < 0.$

Для каждого столбца $X_j$ матрицы $X$ (при $v = e_j$) имеем $u^T X_j \geq 0$, значит $u^T X \in R^d$. Таким образом, выполнено условие (2).

Если же выполнены оба условия, то

$0 > \langle u, y \rangle = u^T (Xv) = (u^T X) v \geq 0,$

так как $u^T X \geq 0$ и $v \geq 0$. Противоречие. Значит, условия не могут выполняться одновременно. Что и требовалось доказать
\end{document}

